% preamble-exam-zh.tex — 共享 preamble
% 用于所有 exam-zh 模板
%
% 样式策略：
%   屏幕 → 语义彩色（蓝=关键想法，红=易错，绿=路线，灰=提示/教师备注）
%   打印 → 自动降级为黑白（通过 print mode 或 monochrome 选项）

\documentclass{exam-zh}

% ---------- 基础配置 ----------
\examsetup{
  page/size = a4paper,
  page/show-head = false,
  page/show-foot = false,
  sealline/show = false,
  font = times,
  math-font = xits,
}

\usepackage{tcolorbox}
\usepackage{needspace}
\usepackage{graphicx}
\usepackage{tikz}
\usetikzlibrary{calc,intersections,angles,quotes,arrows.meta,decorations.markings,patterns.meta,backgrounds}
\usepackage{pgfplots}
\pgfplotsset{compat=1.18}
\tcbuselibrary{skins,breakable}

% ---------- 打印/屏幕颜色切换 ----------
% 用法：\documentclass[monochrome]{exam-zh} 即可切换为纯黑白
% 注意：不要使用 \if@xxx 放在 makeatother 外部；这里改成普通条件名。
\newif\ifeduprintcolor
\eduprintcolortrue

\makeatletter
\@ifclasswith{exam-zh}{monochrome}{\eduprintcolorfalse}{}
\makeatother

% ---------- 语义颜色定义 ----------
% 屏幕：有色彩区分；打印：降级为灰度
\ifeduprintcolor
  % --- 屏幕彩色模式 ---
  \definecolor{edu-blue}{RGB}{47, 82, 122}       % 关键想法
  \definecolor{edu-blue-bg}{RGB}{243, 246, 252}
  \definecolor{edu-red}{RGB}{180, 40, 40}         % 易错提醒
  \definecolor{edu-red-bg}{RGB}{255, 245, 240}
  \definecolor{edu-green}{RGB}{34, 100, 60}       % 路线图
  \definecolor{edu-green-bg}{RGB}{242, 250, 245}
  \definecolor{edu-orange}{RGB}{160, 90, 0}       % 训练目标
  \definecolor{edu-orange-bg}{RGB}{255, 248, 235}
  \definecolor{edu-gray}{RGB}{100, 100, 100}      % 提示/教师备注
  \definecolor{edu-gray-bg}{RGB}{248, 248, 248}
  \definecolor{edu-purple}{RGB}{90, 50, 120}      % 思考题
  \definecolor{edu-purple-bg}{RGB}{248, 244, 252}
\else
  % --- 打印黑白模式 ---
  \definecolor{edu-blue}{gray}{0.15}
  \definecolor{edu-blue-bg}{gray}{0.96}
  \definecolor{edu-red}{gray}{0.25}
  \definecolor{edu-red-bg}{gray}{0.97}
  \definecolor{edu-green}{gray}{0.20}
  \definecolor{edu-green-bg}{gray}{0.97}
  \definecolor{edu-orange}{gray}{0.25}
  \definecolor{edu-orange-bg}{gray}{0.98}
  \definecolor{edu-gray}{gray}{0.40}
  \definecolor{edu-gray-bg}{gray}{0.97}
  \definecolor{edu-purple}{gray}{0.30}
  \definecolor{edu-purple-bg}{gray}{0.98}
\fi
\colorlet{edu-diagram-muted}{edu-gray}

% ---------- TikZ diagram semantic macros ----------
\definecolor{edu-diagram-ink}{RGB}{17,24,39}
\definecolor{edu-diagram-marker}{RGB}{220,38,38}
\definecolor{edu-diagram-angle}{RGB}{5,150,105}
\tikzset{
  diagram point/.style={fill=edu-diagram-ink},
  diagram tick/.style={draw=edu-diagram-marker,line width=0.9pt,line cap=round},
  diagram segment/.style={draw=edu-diagram-ink,line cap=round},
  point label/.style={inner sep=1pt},
  condition label/.style={inner sep=1pt}
}
\providecommand{\DrawSegment}[3][]{\draw[diagram segment,#1] (#2) -- (#3);}
\providecommand{\DrawDashedSegment}[3][]{\draw[diagram segment,dashed,#1] (#2) -- (#3);}
\providecommand{\Triangle}[4][]{\path[#1] (#2) -- (#3) -- (#4) -- cycle;}
\providecommand{\Quadrilateral}[5][]{\path[#1] (#2) -- (#3) -- (#4) -- (#5) -- cycle;}
\providecommand{\PolygonPath}[2][]{\path[#1] #2 -- cycle;}
\providecommand{\DiagramPointRadius}{0.052cm}
\providecommand{\PointDot}[2][]{\fill[diagram point,#1] (#2) circle[radius=\DiagramPointRadius];}
\providecommand{\PointLabel}[3][]{\node[point label,#1] at (#2) {#3};}
\providecommand{\SegmentLabel}[4][]{\node[condition label,#1] at ($(#2)!0.5!(#3)$) {#4};}
\providecommand{\CoordinateTag}[3][]{\node[condition label,#1] at (#2) {#3};}
\providecommand{\AngleMark}[4][]{\pic[draw=edu-diagram-angle,line width=0.9pt,angle radius=0.48cm,#1] {angle=#2--#3--#4};}
\providecommand{\AngleLabel}[5][]{\pic["{#5}",draw=none,angle radius=0.64cm,#1] {angle=#2--#3--#4};}
\providecommand{\RightAngleMark}[4][]{\pic[draw=edu-diagram-marker,line width=0.9pt,angle radius=0.28cm,#1] {right angle=#2--#3--#4};}
\providecommand{\EqualTick}[3][]{%
  \path[postaction={decorate},decoration={markings,mark=at position .5 with {\draw[diagram tick,#1] (0pt,-4pt) -- (0pt,4pt);}}] (#2) -- (#3);%
}
\providecommand{\DoubleEqualTick}[3][]{%
  \path[postaction={decorate},decoration={markings,mark=at position .46 with {\draw[diagram tick,#1] (0pt,-4pt) -- (0pt,4pt);},mark=at position .54 with {\draw[diagram tick,#1] (0pt,-4pt) -- (0pt,4pt);}}] (#2) -- (#3);%
}
\providecommand{\TripleEqualTick}[3][]{%
  \path[postaction={decorate},decoration={markings,mark=at position .43 with {\draw[diagram tick,#1] (0pt,-4pt) -- (0pt,4pt);},mark=at position .5 with {\draw[diagram tick,#1] (0pt,-4pt) -- (0pt,4pt);},mark=at position .57 with {\draw[diagram tick,#1] (0pt,-4pt) -- (0pt,4pt);}}] (#2) -- (#3);%
}
\providecommand{\ParallelMark}[3][]{\EqualTick[#1]{#2}{#3}}
\providecommand{\NamedSegmentPath}[3]{\path[name path=#1] (#2) -- (#3);}
\providecommand{\IntersectPaths}[3]{\path[name intersections={of=#2 and #3, by=#1}];}

% ---------- 自定义教学环境 ----------
% 共性：enhanced + breakable（跨页不断裂）

% 原题卡片 — 强边框，突出显示
\newtcolorbox{problemcard}{
  enhanced, breakable,
  colback=edu-blue-bg,
  colframe=edu-blue,
  boxrule=1.2pt,
  arc=0pt,
  left=3mm, right=3mm, top=3mm, bottom=3mm,
}

% 解题路线图 — 左边线 + 浅底
\newtcolorbox{routebox}{
  enhanced, breakable,
  colback=edu-green-bg,
  colframe=edu-green!30,
  boxrule=0pt,
  borderline west={2.5pt}{0pt}{edu-green},
  left=3mm, right=3mm, top=3mm, bottom=3mm,
}

% 关键想法 — 蓝色左边线
\newtcolorbox{keyideabox}{
  enhanced, breakable,
  colback=edu-blue-bg,
  colframe=edu-blue!30,
  boxrule=0pt,
  borderline west={2.5pt}{0pt}{edu-blue},
  left=3mm, right=3mm, top=3mm, bottom=3mm,
}

% 易错提醒 — 红色左边线
\newtcolorbox{mistakebox}{
  enhanced, breakable,
  colback=edu-red-bg,
  colframe=edu-red!20,
  boxrule=0pt,
  borderline west={3pt}{0pt}{edu-red},
  left=3mm, right=3mm, top=2.5mm, bottom=2.5mm,
}

% 提示 — 灰色虚线左边线
\newtcolorbox{hintbox}{
  enhanced, breakable,
  colback=edu-gray-bg,
  colframe=edu-gray!20,
  boxrule=0pt,
  borderline west={2pt}{0pt}{edu-gray, dashed},
  left=3mm, right=3mm, top=2mm, bottom=2mm,
}

% 思考题 — 紫色虚线左边线
\newtcolorbox{thinkbox}{
  enhanced, breakable,
  colback=edu-purple-bg,
  colframe=edu-purple!20,
  boxrule=0pt,
  borderline west={2pt}{0pt}{edu-purple, dashed},
  left=2.5mm, right=2.5mm, top=2.5mm, bottom=2.5mm,
}

% 教师备注 — 灰色左边线，小字
\newtcolorbox{teachernote}{
  enhanced, breakable,
  colback=edu-gray-bg,
  colframe=edu-gray!15,
  boxrule=0pt,
  borderline west={2pt}{0pt}{edu-gray!60},
  left=2.5mm, right=2.5mm, top=2.5mm, bottom=2.5mm,
  fontupper=\small,
  colupper=black!60,
}

% 训练目标 — 橙色左边线，小字
\newtcolorbox{traininggoal}{
  enhanced, breakable,
  colback=edu-orange-bg,
  colframe=edu-orange!15,
  boxrule=0pt,
  borderline west={1.5pt}{0pt}{edu-orange},
  left=2mm, right=2mm, top=1mm, bottom=1mm,
  fontupper=\small,
}

% 答题区 — 无边框，纯留白
\newcommand{\answerarea}[1][40mm]{%
  \vspace{2mm}%
  \begin{tcolorbox}[
    enhanced, breakable,
    colback=white,
    colframe=white,
    boxrule=0pt,
    height=#1,
    valign=top,
    bottom=0mm,
    after skip=0mm,
  ]
  \end{tcolorbox}%
}

\newsavebox{\diagramtikzbox}

\newenvironment{diagramcoltikz}[2]{%
  \begin{minipage}[t]{#1}
    \vspace{0pt}\centering
    \def\diagramcaption{#2}%
    \begin{lrbox}{\diagramtikzbox}%
}{%
    \end{lrbox}%
    \usebox{\diagramtikzbox}%
    \ifx\diagramcaption\empty\else
      \par\vspace{1mm}{\scriptsize\color{edu-diagram-muted}\diagramcaption}
    \fi
  \end{minipage}%
}

\newenvironment{diagramrowtikz}[3]{%
  \begin{minipage}[t]{#1}
    \vspace{0pt}\centering
    \def\diagramlabel{#2}%
    \def\diagramcaption{#3}%
    \begin{lrbox}{\diagramtikzbox}%
}{%
    \end{lrbox}%
    \usebox{\diagramtikzbox}%
    \ifx\diagramlabel\empty\else
      \par\vspace{1mm}{\scriptsize\bfseries\diagramlabel}
    \fi
    \ifx\diagramcaption\empty\else
      \par\vspace{0.5mm}{\scriptsize\color{edu-diagram-muted}\diagramcaption}
    \fi
  \end{minipage}%
}

\newenvironment{answerareawithdiagramtikz}[3]{%
  \vspace{2mm}%
  \begin{tcolorbox}[
    enhanced, breakable,
    colback=white,
    colframe=white,
    boxrule=0pt,
    height=#1,
    valign=top,
    bottom=0mm,
    after skip=0mm,
  ]
  \begin{minipage}[t]{\dimexpr\linewidth-#2-6mm\relax}
    \vspace{0pt}
  \end{minipage}\hfill
  \begin{diagramcoltikz}{#2}{#3}
}{%
  \end{diagramcoltikz}
  \end{tcolorbox}%
}

% ---------- 字体设置 ----------
% exam-zh (ctexbook) already sets CJK fonts via ctex.
% Only override if specific fonts are needed.
% macOS: Songti SC, STHeiti, STFangsong
% Linux: SimSun, SimHei, FangSong

% ---------- 页面设置 ----------
% exam-zh (ctexbook) already loads geometry.
% Override margins if needed:
% \geometry{top=16mm, bottom=16mm, left=14mm, right=14mm}

\examsetup{
  question/show-answer = true,
  fillin/show-answer = true,
  solution/show-solution = show-stay,
}

% ---------- 教师版解答样式 ----------
\newtcolorbox{solutionblock}{
  enhanced, breakable,
  colback=edu-blue-bg,
  colframe=edu-blue!25,
  boxrule=0pt,
  borderline west={2pt}{0pt}{edu-blue!50},
  arc=0pt,
  left=3mm, right=3mm, top=2mm, bottom=2mm,
  before skip=2mm,
  after skip=3mm,
  fontupper=\small,
}

% 解答步骤编号
\newcounter{solstep}
\newcommand{\solstep}[2]{%
  \stepcounter{solstep}%
  \par\medskip
  \noindent\hangindent=6mm
  {\color{edu-blue}\bfseries\textcircled{\scriptsize\thesolstep}\enspace #1}\enspace #2\par
}

\begin{document}

\section*{矩形存在性 Typed Fusion 作业}

\section*{一、链式入口与关键输出（每题 8 分，共 24 分）}


\begin{question}[points=8]
  已知反比例函数 $y=\dfrac{k}{x}$ 的图像在第一象限与坐标轴围成的矩形面积为 $8$，直线 $y=x+2$ 与该反比例函数交于两点。取第一象限交点 $P$，则 $P$ 的坐标与以 $P$ 为顶点的坐标轴矩形面积分别为（\quad）。

  \begin{choices}
    \item $P(2,4)$，面积 $8$
    \item $P(4,2)$，面积 $8$
    \item $P(2,4)$，面积 $4$
    \item $P(-4,-2)$，面积 $8$
  \end{choices}
  \paren[A]
  \begin{solutionblock}
    面积 $8$ 且第一象限，得 $k=8$。解 $x+2=\dfrac8x$，得交点 $(2,4)$ 与 $(-4,-2)$，第一象限交点为 $P(2,4)$，坐标轴矩形面积为 $8$。
  \end{solutionblock}
\end{question}



\begin{question}[points=8]
  反比例函数 $y=\dfrac{k}{x}$ 在第一象限对应的坐标轴矩形面积为 $12$。直线 $y=x+1$ 与它的第一象限交点为 $P$。若以 $A(0,0)$，$B(3,0)$，$C=P$ 为相邻顶点作矩形 $ABCD$，则 $D=$ \underline{\hspace{2cm}}。

  \paren[$(0,4)$]
  \begin{solutionblock}
    面积给出 $k=12$。解 $x+1=\dfrac{12}{x}$，第一象限交点为 $P(3,4)$。固定顺序矩形中 $D=A+C-B=(0,4)$。
  \end{solutionblock}
\end{question}



\begin{question}[points=8]
  已知 $A(0,0)$，$B(6,0)$，点 $C$ 在直线 $y=x$ 上。若 $A,B,C,D$ 构成矩形，且矩形面积为 $18$，再令 $C$ 落在反比例函数 $y=\dfrac{k}{x}$ 上，则 $k=$ \underline{\hspace{2cm}}。

  \paren[$9$]
  \begin{solutionblock}
    构成矩形候选中，面积为 $18$ 的是 $C(3,3),D(3,-3)$。再由 $C$ 在 $y=\dfrac{k}{x}$ 上得 $k=3\cdot3=9$。
  \end{solutionblock}
\end{question}


\section*{二、两段 relation chain（每题 12 分，共 36 分）}

\needspace{8\baselineskip}

\begin{problem}[points=12]
  反比例函数 $y=\dfrac{k}{x}$ 在第一象限对应的坐标轴矩形面积为 $20$。点 $P$ 是该反比例函数上横坐标为 $4$ 的点，直线 $OP$ 与点 $A(0,0)$，$B(4,0)$ 共同确定矩形 $ABCD$，其中 $C=P$。求直线 $OP$ 的解析式、$k$ 的值以及点 $C,D$。

  \answerarea[55mm]
  \setcounter{solstep}{0}
  \begin{solutionblock}
    \textbf{答案：}$k=20$，$C(4,5)$，$D(0,5)$，$OP:y=\dfrac54x$\par\medskip
    \solstep{面积输出反比例函数}{第一象限坐标轴矩形面积为 $20$，所以 $k=20$，反比例函数为 $y=\dfrac{20}{x}$。}
    \solstep{函数输出点}{当 $x_P=4$ 时，$y_P=5$，所以 $P(4,5)$。}
    \solstep{点输出直线和矩形}{直线 $OP$ 为 $y=\dfrac54x$。又 $C=P$，所以 $C(4,5)$，$D=A+C-B=(0,5)$。}
  \end{solutionblock}
\end{problem}


\needspace{8\baselineskip}

\begin{problem}[points=12]
  一次函数 $y=x+2$ 与反比例函数 $y=\dfrac{8}{x}$ 相交。先求两图像交点并得到不等式 $x+2\ge\dfrac{8}{x}$ 的解集；再在直线 $y=x+2$ 上寻找点 $C$，使 $A(0,0)$，$B(2,0)$，$C,D$ 构成矩形，且 $x_C$ 属于上述解集。求所有可能的 $C,D$。

  \answerarea[65mm]
  \setcounter{solstep}{0}
  \begin{solutionblock}
    \textbf{答案：}$C(2,4)$，$D(0,4)$\par\medskip
    \solstep{函数图像输出位置筛选}{由 $x+2=\dfrac8x$ 得交点横坐标 $x=-4,2$。分段比较得 $x+2\ge\dfrac8x$ 的解集为 $\left[-4,0\right)\cup\left[2,+\infty\right)$。}
    \solstep{矩形候选}{设 $C(t,t+2)$。若直角在 $A$，得 $t=0$，但 $0$ 不在筛选区间；若直角在 $B$，得 $t=2$，所以 $C(2,4)$，$D=A+C-B=(0,4)$；直角在 $C$ 无实数解。}
  \end{solutionblock}
\end{problem}


\needspace{8\baselineskip}

\begin{problem}[points=12]
  已知 $A(0,0)$，$B(4,0)$，点 $C$ 在直线 $y=x$ 上。先求使四边形 $ABCD$ 为菱形的 $C,D$；再判断这个候选四边形是否也是矩形。

  \answerarea[55mm]
  \setcounter{solstep}{0}
  \begin{solutionblock}
    \textbf{答案：}$C(4,4)$，$D(0,4)$；也是矩形\par\medskip
    \solstep{菱形候选}{设 $C(t,t)$。固定顺序菱形用 $AB=BC$，列 $(t-4)^2+t^2=16$，得 $t=0$ 或 $t=4$。舍去退化解，得 $C(4,4)$，$D=A+C-B=(0,4)$。}
    \solstep{矩形验证}{$AB$ 水平，$BC$ 竖直，所以 $AB\perp BC$。该候选四边形也是矩形。}
  \end{solutionblock}
\end{problem}


\section*{三、三段 typed chain 综合（每题 10 分，共 40 分）}

\needspace{8\baselineskip}

\begin{problem}[points=10]
  已知 $A(0,0)$，$B(8,0)$。四边形 $ABCD$ 是矩形且面积为 $16$，点 $C$ 在第一象限。设 $C$ 同时在直线 $y=mx$ 与反比例函数 $y=\dfrac{k}{x}$ 上。求 $C,D,m,k$。

  \answerarea[58mm]
  \setcounter{solstep}{0}
  \begin{solutionblock}
    \textbf{答案：}$C(8,2)$，$D(0,2)$，$m=\dfrac14$，$k=16$\par\medskip
    \solstep{矩形面积输出点}{$AB=8$，面积为 $16$，且 $C$ 在第一象限，所以 $BC=2$，得 $C(8,2)$，$D(0,2)$。}
    \solstep{点输出两个函数参数}{$C$ 在 $y=mx$ 上，得 $m=\dfrac{2}{8}=\dfrac14$；在 $y=\dfrac{k}{x}$ 上，得 $k=8\times2=16$。}
  \end{solutionblock}
\end{problem}


\needspace{8\baselineskip}

\begin{problem}[points=10]
  已知 $A(0,0)$，$B(4,0)$，点 $C$ 在直线 $y=2x$ 上。先由矩形 $ABCD$ 生成候选点对，再把该候选点对送入菱形条件检验。判断是否存在同一组 $C,D$，使四边形 $ABCD$ 同时为矩形和菱形。

  \answerarea[52mm]
  \setcounter{solstep}{0}
  \begin{solutionblock}
    \textbf{答案：}不存在\par\medskip
    \solstep{矩形候选}{固定顺序矩形中 $BC\perp AB$，所以 $x_C=4$。又 $C$ 在 $y=2x$ 上，得 $C(4,8)$，$D(0,8)$。}
    \solstep{菱形验证}{$AB=4$，$BC=8$，不满足邻边相等。因此不存在同一组 $C,D$ 同时满足矩形和菱形。}
  \end{solutionblock}
\end{problem}


\needspace{8\baselineskip}

\begin{problem}[points=10]
  反比例函数 $y=\dfrac{k}{x}$ 在第一象限对应的坐标轴矩形面积为 $18$。直线 $y=x+3$ 与该反比例函数交于两点，取第一象限交点 $P$。过 $O(0,0)$ 与 $P$ 作直线 $OP$，再在直线 $OP$ 上取点 $E$，使四边形 $ABEF$ 是菱形，其中 $A(0,0)$，$B(6,0)$。求 $P$、直线 $OP$ 的解析式以及 $E,F$。

  \answerarea[75mm]
  \setcounter{solstep}{0}
  \begin{solutionblock}
    \textbf{答案：}$P(3,6)$，$OP:y=2x$；$E\left(\dfrac{12}{5},\dfrac{24}{5}\right)$，$F\left(-\dfrac{18}{5},\dfrac{24}{5}\right)$\par\medskip
    \solstep{面积输出反比例函数}{第一象限面积为 $18$，得 $k=18$，反比例函数为 $y=\dfrac{18}{x}$。}
    \solstep{交点输出P}{解 $x+3=\dfrac{18}{x}$，得 $x=3$ 或 $x=-6$。第一象限点为 $P(3,6)$。}
    \solstep{P输出直线}{直线 $OP$ 为 $y=2x$。}
    \solstep{直线输出菱形候选}{设 $E(t,2t)$。菱形 $ABEF$ 中 $AB=BE$，列 $(t-6)^2+(2t)^2=36$，得 $t=0$ 或 $t=\dfrac{12}{5}$。舍去退化解，得 $E\left(\dfrac{12}{5},\dfrac{24}{5}\right)$，$F=A+E-B=\left(-\dfrac{18}{5},\dfrac{24}{5}\right)$。}
  \end{solutionblock}
\end{problem}


\needspace{8\baselineskip}

\begin{problem}[points=10]
  反比例函数 $y=\dfrac{k}{x}$ 在第一象限对应的坐标轴矩形面积为 $12$。直线 $y=x+1$ 与该反比例函数交于两点，取第一象限交点 $P$。以 $A(0,0)$，$B(x_P,0)$，$C=P$ 构成矩形 $ABCD$；再用点 $C$ 反求一次函数 $y=mx$。求 $k,P,D,m$，并说明每一步的输出进入了下一步的哪个输入。

  \answerarea[75mm]
  \setcounter{solstep}{0}
  \begin{solutionblock}
    \textbf{答案：}$k=12$，$P(3,4)$，$D(0,4)$，$m=\dfrac43$\par\medskip
    \solstep{面积输出反比例函数}{第一象限面积为 $12$，得 $k=12$。}
    \solstep{函数输出交点}{解 $x+1=\dfrac{12}{x}$，得 $x=3$ 或 $x=-4$，第一象限点 $P(3,4)$。}
    \solstep{点输出矩形}{$B(x_P,0)=B(3,0)$，$C=P(3,4)$，所以 $D=A+C-B=(0,4)$。}
    \solstep{点输出一次函数参数}{$C$ 在 $y=mx$ 上，所以 $m=\dfrac43$。}
  \end{solutionblock}
\end{problem}


\clearpage
\section*{答案速查}




\end{document}
